\documentclass[conference]{IEEEtran}
\usepackage{cite}
\usepackage{url}
\usepackage{hyperref}
\usepackage{booktabs}
\usepackage{multirow}
\usepackage{array}
\usepackage{makecell}
\usepackage{amsmath}
\usepackage{graphicx}

\begin{document}

\title{Supplementary Material:\\``The Cloud EC2 Sweet Spot: Cost Optimization and Scaling Strategies\\for Video-as-a-Service Workloads''}

\maketitle

\tableofcontents

\bigskip\noindent
All experimental materials are available at the artifact repository:\\
\url{https://anonymous.4open.science/r/ec2sweetspot-blinded2}

% -----------------------------------------------------------------------
\section{EC2 Instance Hardware Specifications}
\label{app:specs}
% -----------------------------------------------------------------------

Table~\ref{tab:app_instance_specs} documents the hardware configuration and
on-demand Linux pricing for all 17 EC2 instance types evaluated in this study.
Prices reflect the AWS Pricing API for \texttt{us-east-1} as of March~2026.

\begin{table*}[!htbp]
\centering
\caption{EC2 instance specifications and on-demand pricing (\texttt{us-east-1},
Linux, March~2026). T-family instances operate in Unlimited burst mode; sustained
100\% vCPU utilization beyond the credit baseline may incur surplus credit
surcharges (see main paper, Threats to Validity).}
\label{tab:app_instance_specs}
\small
\begin{tabular}{llrrll}
\toprule
\textbf{Instance} & \textbf{Family} & \textbf{vCPU} & \textbf{RAM (GiB)}
  & \textbf{Dominant Processor (Phase~1)} & \textbf{Price (\$/hr)} \\
\midrule
\texttt{t2.micro}    & Burstable (T) & 1  & 1   & Intel Xeon E5-2686 v4 / Platinum 8259CL & 0.0116 \\
\texttt{t3.micro}    & Burstable (T) & 2  & 1   & Intel Xeon Platinum 8259CL (2.50\,GHz)  & 0.0104 \\
\texttt{t3.small}    & Burstable (T) & 2  & 2   & Intel Xeon Platinum 8259CL (2.50\,GHz)  & 0.0208 \\
\texttt{t3.medium}   & Burstable (T) & 2  & 4   & Intel Xeon Platinum 8259CL (2.50\,GHz)  & 0.0416 \\
\texttt{t3.large}    & Burstable (T) & 2  & 8   & Intel Xeon Platinum 8259CL (2.50\,GHz)  & 0.0832 \\
\texttt{t3.xlarge}   & Burstable (T) & 4  & 16  & Intel Xeon Platinum 8259CL (2.50\,GHz)  & 0.1664 \\
\texttt{t3.2xlarge}  & Burstable (T) & 8  & 32  & Intel Xeon Platinum 8259CL (2.50\,GHz)  & 0.3328 \\
\midrule
\texttt{m5.large}    & General (M)   & 2  & 8   & Intel Xeon Platinum 8259CL (2.50\,GHz)  & 0.0960 \\
\texttt{m5.xlarge}   & General (M)   & 4  & 16  & Intel Xeon Platinum 8259CL (2.50\,GHz)  & 0.1920 \\
\texttt{m5.2xlarge}  & General (M)   & 8  & 32  & Intel Xeon Platinum 8259CL (2.50\,GHz)  & 0.3840 \\
\texttt{m5.4xlarge}  & General (M)   & 16 & 64  & Intel Xeon Platinum 8259CL (2.50\,GHz)  & 0.7680 \\
\midrule
\texttt{c5.large}    & Compute (C)   & 2  & 4   & Intel Xeon Platinum 8275CL (3.00\,GHz)  & 0.0850 \\
\texttt{c5.xlarge}   & Compute (C)   & 4  & 8   & Intel Xeon Platinum 8275CL (3.00\,GHz)  & 0.1700 \\
\texttt{c5.2xlarge}  & Compute (C)   & 8  & 16  & Intel Xeon Platinum 8275CL (3.00\,GHz)  & 0.3400 \\
\texttt{c5.4xlarge}  & Compute (C)   & 16 & 32  & Intel Xeon Platinum 8275CL (3.00\,GHz)  & 0.6800 \\
\midrule
\texttt{m7g.large}   & ARM Graviton~3 & 2 & 8  & AWS Graviton~3 (Arm64, 100\% uniform)   & 0.0808 \\
\texttt{c7g.large}   & ARM Graviton~3 & 2 & 4  & AWS Graviton~3 (Arm64, 100\% uniform)   & 0.0725 \\
\bottomrule
\end{tabular}
\end{table*}

% -----------------------------------------------------------------------
\section{Phase 1: CPU Homogeneity Validation}
\label{app:phase1}
% -----------------------------------------------------------------------

\subsection{Protocol}

For each of the 15 x86 instance types, approximately 100 instances were deployed
across all six \texttt{us-east-1} availability zones (\texttt{us-east-1a}
through \texttt{us-east-1f}). Immediately upon SSH readiness, each instance's
processor model was read from \texttt{/proc/cpuinfo}. Three Intel Xeon
generations were observed in the x86 fleet: Xeon~E5-2686~v4 (Broadwell),
Xeon~8175M (Skylake-SP), and Xeon~Platinum~8259CL / 8275CL (Cascade Lake).
ARM Graviton~3 instances (\texttt{m7g.large}, \texttt{c7g.large}) exhibited
complete processor uniformity across all sampled nodes.

\subsection{Results and Hardware Standardization Policy}

Table~\ref{tab:app_cpu_homogeneity} summarizes the dominant processor and
minimum homogeneity rate observed per instance family. The M-family exhibited
the greatest heterogeneity (60/40 split between Platinum~8259CL and Xeon~8175M
in zone~\texttt{c}), and the C-family showed a 74/26 split between
Platinum~8275CL and Xeon~8124M in zone~\texttt{b}. T-family instances were
nearly uniform (98\% 8259CL in zone~\texttt{a}).

\begin{table}[!htbp]
\centering
\caption{Phase~1 CPU homogeneity results. \emph{Min.\ Dominant \%} is the
lowest fraction of spawned instances matching the target processor across all
sampled zones. All Phase~2 and Phase~3 benchmarks were executed exclusively on
the dominant processor via the terminate-and-retry protocol.}
\label{tab:app_cpu_homogeneity}
\small
\begin{tabular}{lllr}
\toprule
\textbf{Family} & \textbf{Target CPU} & \textbf{Clock} & \textbf{Min.\ Dominant \%} \\
\midrule
T (\texttt{t3.*}) & Xeon Platinum 8259CL & 2.50\,GHz & 98\% \\
M (\texttt{m5.*}) & Xeon Platinum 8259CL & 2.50\,GHz & 60\% \\
C (\texttt{c5.*}) & Xeon Platinum 8275CL & 3.00\,GHz & 74\% \\
ARM (\texttt{m7g.large}) & AWS Graviton~3      & ---       & 100\% \\
ARM (\texttt{c7g.large}) & AWS Graviton~3      & ---       & 100\% \\
\bottomrule
\end{tabular}
\end{table}

To eliminate hardware heterogeneity as a confounding variable, all Phase~2 and
Phase~3 benchmarks enforced the target processor at runtime: each provisioned
instance was verified via \texttt{/proc/cpuinfo} before workload dispatch, and
non-conforming instances were immediately terminated and replaced. This terminate-and-retry loop introduced a per-instance verification overhead
of less than 2\,s, ensuring that all 30 runs per configuration were executed
on hardware-homogeneous nodes.

% -----------------------------------------------------------------------
\section{Phase 2: Statistical Analysis Details}
\label{app:stats}
% -----------------------------------------------------------------------

\subsection{Normality Assessment}

For each of the $15 \times 4 \times 3 = 180$ Phase~2 configurations, a
Shapiro-Wilk test ($\alpha = 0.05$, $n = 30$) was applied to assess distributional
assumptions. The majority of configurations did not reject normality.
Non-normal cases arose primarily for single-thread configurations (elevated
variability in FFmpeg process startup and thread scheduling jitter) and for VP9
on larger instances (multimodal timing distributions associated with thread
affinity scheduling behavior). Where normality was rejected, the Central Limit
Theorem ($n = 30$) justifies retention of Welch's t-test; Mann-Whitney~U was
applied as a non-parametric confirmatory test where results were ambiguous.
Complete Shapiro-Wilk statistics per configuration are available in
\texttt{analysis\_output/statistical\_tests\_c5.csv} and
\texttt{analysis\_output/statistical\_tests\_m5.csv} in the artifact repository.

\subsection{Phase 3 ARM Graviton: Key Inferential Tests}

Two inferential tests were applied to ARM Graviton Phase~3 results, as
reported in the main paper:

\textbf{Mann-Whitney~U (\texttt{m7g.large} vs.\ \texttt{c7g.large}, H.265
horizontal):} $U = 432$, $p = 0.71$. The null hypothesis of equal medians is
not rejected. The high coefficient of variation (CV~$>$50\%, driven by bimodal
provisioning outcomes) limits statistical power for this test.

\textbf{Welch's t-test (\texttt{m7g.large} H.264, serial vs.\ horizontal):}
$t(30.2) = -14.7$, $p < 10^{-12}$, Cohen's $d = 4.31$. The very large effect
size confirms that the H.264 performance inversion is driven by provisioning
overhead ($\approx$48\,s) substantially exceeding the per-chunk encoding
duration ($\approx$6\,s) for this configuration.

% -----------------------------------------------------------------------
\section{Phase 3: Provisioning Overhead Decomposition}
\label{app:overhead}
% -----------------------------------------------------------------------

Table~\ref{tab:app_overhead} provides the complete provisioning overhead
percentages for all Phase~3 configurations. Setup overhead is defined as
$(T_{\text{total}} - T_{\text{enc}}) / T_{\text{total}}$, where $T_{\text{enc}}$
is the active FFmpeg encoding duration. These values were omitted from
Table~V of the main paper due to space constraints.

\begin{table*}[!htbp]
\centering
\caption{Provisioning overhead decomposition for all Phase~3 configurations.
\textbf{Dynamic} (Phase~3.1): standard Ubuntu AMI with FFmpeg installed via
\texttt{apt} at runtime. \textbf{Optimized} (Phase~3.2): custom pre-built AMI
with FFmpeg pre-installed. Setup OH\% $= (T_{\mathrm{total}} -
T_{\mathrm{enc}}) / T_{\mathrm{total}}$, where $T_{\mathrm{enc}}$ is the
active FFmpeg encoding duration on the critical path. All times are median
wall-clock ($n=30$, $\pm$\,1.96\,$\hat{\sigma}/\sqrt{n}$).}
\label{tab:app_overhead}
\small
\begin{tabular}{llcc}
\toprule
\textbf{Instance / Strategy} & \textbf{Provisioning} & \textbf{Total Time (min)} & \textbf{Setup OH} \\
\midrule
\multicolumn{4}{c}{\textit{Burstable T-Family}} \\
\midrule
\texttt{t3.micro} serial          & Dynamic   & 4.03 $\pm$ 0.12 & 57.3\% \\
\texttt{t3.micro} serial          & Optimized & 2.72 $\pm$ 0.08 & 36.8\% \\
10$\times$\texttt{t3.micro} horiz.& Dynamic   & 3.28 $\pm$ 0.10 & 93.3\% \\
10$\times$\texttt{t3.micro} horiz.& Optimized & 1.22 $\pm$ 0.04 & 82.0\% \\
\texttt{t3.2xlarge} vertical      & Dynamic   & 3.04 $\pm$ 0.09 & 74.0\% \\
\texttt{t3.2xlarge} vertical      & Optimized & 1.80 $\pm$ 0.05 & 56.1\% \\
\midrule
\multicolumn{4}{c}{\textit{General Purpose M-Family}} \\
\midrule
\texttt{m5.large} serial          & Dynamic   & 3.46 $\pm$ 0.10 & 50.6\% \\
\texttt{m5.large} serial          & Optimized & 2.71 $\pm$ 0.08 & 36.9\% \\
10$\times$\texttt{m5.large} horiz.& Dynamic   & 3.63 $\pm$ 0.11 & 92.8\% \\
10$\times$\texttt{m5.large} horiz.& Optimized & 1.22 $\pm$ 0.04 & 78.7\% \\
\texttt{m5.4xlarge} vertical      & Dynamic   & 2.33 $\pm$ 0.07 & 74.7\% \\
\texttt{m5.4xlarge} vertical      & Optimized & 1.56 $\pm$ 0.05 & 62.2\% \\
\midrule
\multicolumn{4}{c}{\textit{Compute Optimized C-Family}} \\
\midrule
\texttt{c5.large} serial          & Dynamic   & 3.21 $\pm$ 0.10 & 53.9\% \\
\texttt{c5.large} serial          & Optimized & 2.46 $\pm$ 0.07 & 39.8\% \\
10$\times$\texttt{c5.large} horiz.& Dynamic   & 3.52 $\pm$ 0.11 & 93.2\% \\
10$\times$\texttt{c5.large} horiz.& Optimized & 1.18 $\pm$ 0.04 & 79.7\% \\
\texttt{c5.4xlarge} vertical      & Dynamic   & 2.13 $\pm$ 0.06 & 75.6\% \\
\texttt{c5.4xlarge} vertical      & Optimized & 1.49 $\pm$ 0.04 & 65.1\% \\
\bottomrule
\end{tabular}
\end{table*}

The 93.3\% overhead for 10$\times$\texttt{t3.micro} Dynamic (Phase~3.1)
represents the upper bound in this study: ten instances are launched and
configured in parallel, but the synchronization barrier at output merge means
the slowest provisioning event determines wall-clock time. Under the optimized
AMI (Phase~3.2), the overhead for this configuration falls to 82\%, reflecting
the minimum EC2 instance launch latency observed under these conditions
($\approx$48\,s: EBS volume attachment, OS initialization, SSH readiness).
Serial single-instance configurations exhibit lower relative overhead
(37--54\% Dynamic, 37--40\% Optimized) because a single instance is
provisioned before useful work begins.

% -----------------------------------------------------------------------
\section{ARM Graviton Independent Replication}
\label{app:rerun}
% -----------------------------------------------------------------------

To independently validate the ARM Graviton results reported in Table~III of the
main paper, all eight ARM configurations were replicated using a hardware
enforcement protocol consistent with the terminate-and-retry methodology
established in Phase~1. The validation script (\texttt{rerun\_cpu\_verified.py})
verifies each provisioned instance's processor model via \texttt{/proc/cpuinfo}
before workload dispatch, immediately terminating and replacing any
non-conforming node.

Table~\ref{tab:app_rerun_summary} reports the replication completion status.
All 240 runs (8 configurations $\times$ 30) confirmed the expected Graviton~3
processor prior to encoding. Performance results from this independent
replication were statistically indistinguishable from the primary dataset,
confirming the reproducibility of the ARM Graviton findings. Full replication
logs are available in \texttt{phase3\_scaling/n30\_rerun\_logs/} in the
artifact repository.

\begin{table}[!htbp]
\centering
\caption{ARM Graviton independent replication summary.
All 8 configurations completed $n=30$ runs with hardware-verified processors.
Replication logs available at \texttt{phase3\_scaling/n30\_rerun\_logs/}.}
\label{tab:app_rerun_summary}
\small
\begin{tabular}{lllc}
\toprule
\textbf{Instance} & \textbf{Strategy} & \textbf{Codec} & \textbf{Runs} \\
\midrule
\texttt{m7g.large} & Serial     & H.264 & 30 \\
\texttt{m7g.large} & Serial     & H.265 & 30 \\
\texttt{m7g.large} & Serial     & VP9   & 30 \\
\texttt{m7g.large} & Horizontal & H.265 & 30 \\
\texttt{c7g.large} & Serial     & H.264 & 30 \\
\texttt{c7g.large} & Serial     & H.265 & 30 \\
\texttt{c7g.large} & Serial     & VP9   & 30 \\
\texttt{c7g.large} & Horizontal & H.264 & 30 \\
\bottomrule
\end{tabular}
\end{table}

% -----------------------------------------------------------------------
\section{Regional Cost Comparison (\texttt{sa-east-1})}
\label{app:regional}
% -----------------------------------------------------------------------

\subsection{Pricing Differential}

Table~\ref{tab:app_regional_pricing} documents the on-demand hourly prices for
the three horizontal scaling instance types in both evaluated regions, obtained
from the AWS Pricing API in March~2026. The observed 54--62\% premium for
\texttt{sa-east-1} is higher than the 43--44\% differential cited in prior
literature, reflecting AWS pricing revisions to the São Paulo region since those
studies were conducted.

\begin{table}[!htbp]
\centering
\caption{On-demand Linux pricing (\$/hr) for horizontal scaling instance types,
\texttt{us-east-1} vs.\ \texttt{sa-east-1}, March~2026. Premium is the
percentage increase of \texttt{sa-east-1} relative to \texttt{us-east-1}.}
\label{tab:app_regional_pricing}
\small
\begin{tabular}{lrrr}
\toprule
\textbf{Instance} & \textbf{us-east-1} & \textbf{sa-east-1} & \textbf{Premium} \\
\midrule
\texttt{t3.micro}  & \$0.01040 & \$0.01677 & $+$61.5\% \\
\texttt{m5.large}  & \$0.09600 & \$0.15302 & $+$59.4\% \\
\texttt{c5.large}  & \$0.08500 & \$0.13099 & $+$54.1\% \\
\bottomrule
\end{tabular}
\end{table}

\subsection{Performance and Cost Results}

Table~\ref{tab:app_regional_results} presents the empirical H.264 horizontal
scaling results for both regions (10$\times$ nodes, 60\,s chunks, Phase~3.2 AMI,
$n=30$ per configuration). Wall-clock time differences between regions are
consistent with AZ-level provisioning latency variability; the underlying CPU
architecture and clock frequency are equivalent across regions for the same
instance type.

\begin{table*}[!htbp]
\centering
\caption{Regional comparison: H.264 horizontal scaling (10$\times$ nodes,
60\,s chunks, Phase~3.2 AMI, $n=30$). Values are median end-to-end
wall-clock time and total cluster cost per run.}
\label{tab:app_regional_results}
\small
\begin{tabular}{lrrrr}
\toprule
\textbf{Instance} & \textbf{\texttt{us-east-1} Time (s)} & \textbf{\texttt{us-east-1} Cost}
  & \textbf{\texttt{sa-east-1} Time (s)} & \textbf{\texttt{sa-east-1} Cost} \\
\midrule
10$\times$\texttt{t3.micro}  & 73.20  & \$0.0021 & 96.64  & \$0.0045 \\
10$\times$\texttt{c5.large}  & 64.04  & \$0.0152 & 64.41  & \$0.0234 \\
10$\times$\texttt{m5.large}  & 56.41  & \$0.0150 & 79.85  & \$0.0339 \\
\bottomrule
\end{tabular}
\end{table*}

The \texttt{c5.large} cluster exhibited nearly identical wall-clock performance
across regions (64.04\,s vs.\ 64.41\,s), with cost difference driven entirely
by the regional pricing differential. The \texttt{t3.micro} cluster showed
greater performance variability in \texttt{sa-east-1} (96.64\,s vs.\ 73.20\,s),
consistent with higher provisioning latency observed in that region. The
\texttt{m5.large} cluster exhibited intermediate variability (79.85\,s vs.\
56.41\,s) with a correspondingly higher cost penalty.

% -----------------------------------------------------------------------
\section{Reproduction Requirements}
\label{app:repro}
% -----------------------------------------------------------------------

\subsection{Infrastructure}

\begin{itemize}
  \item AWS account with EC2 access in \texttt{us-east-1} and \texttt{sa-east-1}
  \item IAM permissions: \texttt{ec2:RunInstances}, \texttt{ec2:TerminateInstances},
        \texttt{ec2:DescribeInstances}, \texttt{ec2:CreateKeyPair},
        \texttt{ec2:AuthorizeSecurityGroupIngress}
  \item EC2 key pair for SSH-based workload dispatch and result retrieval
  \item Support for all 17 instance types documented in
        Table~\ref{tab:app_instance_specs}
  \item Ubuntu~22.04~LTS base AMI (\texttt{us-east-1}:
        \texttt{ami-0a313d6098716f372}; for other regions, use the equivalent
        Ubuntu~22.04~LTS AMI from the AWS EC2 catalog)
  \item Phase~3.2 optimized AMI with pre-installed FFmpeg and Docker
        (\texttt{us-east-1} x86: \texttt{ami-06c7c20c67513469a}; ARM AMI
        must be built from the provided configuration scripts)
\end{itemize}

\subsection{Software Dependencies}

\begin{itemize}
  \item Python~3.8+ with packages: \texttt{boto3}, \texttt{paramiko},
        \texttt{pandas}, \texttt{numpy}, \texttt{matplotlib}, \texttt{scipy}
  \item FFmpeg~4.4+ with \texttt{libx264}, \texttt{libx265}, and
        \texttt{libvpx-vp9} codec support (pre-installed in optimized AMI)
  \item AWS CLI~v2 with configured credentials
  \item Docker Engine (containerized execution environment in Phase~3 runs)
\end{itemize}

\subsection{Benchmark Workload}

The primary benchmark workload is a 10-minute, 720p (1280$\times$720), H.264
source video, 56\,MB on disk. For the primary experiments, this video is
segmented into 10 chunks of 60\,s each using
\texttt{ffmpeg -f segment -segment\_time 60}. The chunk sensitivity analysis
uses 20 chunks (30\,s) and 5 chunks (120\,s). Substitute videos of equivalent
duration and resolution may be used for reproduction; absolute encoding times
will differ, while the relative scaling trends and overhead ratios are
determined by the workload decomposition structure rather than the specific
video content.

\subsection{Estimated Reproduction Cost and Time}

\begin{itemize}
  \item Phase~1 (CPU homogeneity, $\approx$4,100 short-lived instances):
        \$15--\$25
  \item Phase~2 ($15 \times 4 \times 3 \times 30 = 5{,}400$ runs):
        \$50--\$80
  \item Phase~3 x86 scaling ($18$ configurations $\times$ 30 runs):
        \$100--\$150
  \item Phase~3 ARM Graviton (8 configurations $\times$ 30 runs each, two datasets):
        \$80--\$120
  \item Phase~4 chunk sensitivity (9 configurations $\times$ 30 runs):
        \$40--\$60
  \item \texttt{sa-east-1} regional replication (3 configurations $\times$ 30 runs):
        \$60--\$100
  \item AMI build, data transfer, and miscellaneous: \$30--\$60
\end{itemize}

\noindent
\textbf{Total:} \$375--\$595 at AWS on-demand pricing (March~2026).
\textbf{Estimated wall-clock duration:} 4--6 days, including AMI build time
and sequential phase dependencies.

% -----------------------------------------------------------------------
\section{Artifact Repository Structure}
\label{app:structure}
% -----------------------------------------------------------------------

The artifact repository contains the following principal components:

\begin{verbatim}
ec2sweetspot/
+-- artifacts2/
|   +-- phase1_homogeneity/
|   |   +-- cpu_verification_logs/  (4,100+ cpuinfo records)
|   |   +-- analyze_cpu_distribution.py
|   |   +-- detect_cpu_homogeneity.py
|   +-- phase2_benchmarks/
|   |   +-- all_instances_phase2_extended.py
|   +-- phase3_scaling/
|   |   +-- phase3_unified.py       (Dynamic, Phase 3.1)
|   |   +-- phase3_unified_opt.py   (Optimized AMI, Phase 3.2)
|   |   +-- rerun_cpu_verified.py   (ARM independent replication)
|   |   +-- n30_logs/               (Phase 3.1 execution logs)
|   |   +-- n30_rerun_logs/         (ARM replication logs)
|   |   +-- old_logs_backup/        (ARM horizontal logs)
|   +-- experimental_data/
|   |   +-- aggregated_*.csv        (9 aggregated result files)
|   +-- analysis_output/
|   |   +-- statistical_tests_*.csv (Shapiro-Wilk per config)
|   +-- scripts/processing/
|       +-- regenerate_figures.py
|       +-- parse_n30_logs.py
|       +-- generate_arm_metrics.py
+-- cloud2026/
|   +-- cloud2026.tex               (main paper)
|   +-- appendix.tex                (this document)
|   +-- graphs/                     (generated figure files)
+-- test_videos/
|   +-- bolofofos10min.mp4          (720p, 56 MB, 10 min)
+-- CHECKLIST.md
\end{verbatim}

\subsection{Key Scripts}

\textbf{\texttt{phase3\_unified\_opt.py}:}
Orchestrates Phase~3.2 experiments. Launches EC2 instances using the pre-built
AMI (\texttt{ami-06c7c20c67513469a}), verifies the processor model, transfers
video chunks via SFTP, dispatches parallel FFmpeg encoding, retrieves encoded
output, and terminates instances. Implements \texttt{boto3} instance waiters
with a 30\,s SSH readiness buffer.

\textbf{\texttt{rerun\_cpu\_verified.py}:}
Orchestrates ARM Graviton independent replication. Architecturally identical to
\texttt{phase3\_unified\_opt.py} with runtime CPU enforcement: any instance
returning an unexpected model string is terminated and replaced before workload
dispatch.

\textbf{\texttt{regenerate\_figures.py}:}
Reconstructs both paper figures (\texttt{unified\_codecs\_bar\_chart.png} and
\texttt{unified\_scaling\_comparison.png}) from Phase~2 data embedded in the
paper's LaTeX source, using \texttt{matplotlib}. Enables figure reproduction
without access to raw log files.

% -----------------------------------------------------------------------
\section{Result Traceability}
\label{app:validation}
% -----------------------------------------------------------------------

Each quantitative result in the main paper is traceable to a specific artifact:

\begin{itemize}
  \item \textbf{Table~II (Phase~2 benchmarks):}
        \texttt{experimental\_data/aggregated\_*.csv} (one file per instance type,
        containing per-run timing and cost data).
  \item \textbf{Table~III (ARM Graviton):}
        \texttt{phase3\_scaling/m3\_graviton\_results\_consolidated.csv}
        and execution logs in \texttt{old\_logs\_backup/}. All configurations
        with $n=30$ are directly reproducible from the artifact; the
        \texttt{m7g.large} H.264 horizontal entry ($n=24$) corresponds to
        the available completed runs for that configuration.
  \item \textbf{Table~V (Phase~3 scaling):}
        \texttt{phase3\_scaling/n30\_logs/} (Phase~3.1) and
        \texttt{phase3\_scaling/n30\_parsed/averages\_n30.csv} (Phase~3.2
        summary statistics).
  \item \textbf{Table~VI (Chunk sensitivity):}
        Chunk sensitivity experiment logs in \texttt{phase3\_scaling/}.
  \item \textbf{Regional results (Section~V-C):}
        \texttt{phase3\_scaling/m2\_results\_sa\_east\_1.csv}.
  \item \textbf{Statistical tests (Section~V-D):}
        \texttt{analysis\_output/statistical\_tests\_c5.csv} and
        \texttt{analysis\_output/statistical\_tests\_m5.csv}.
\end{itemize}

\subsection{Reproducibility Checklist}

The \texttt{CHECKLIST.md} file in the repository guides independent
reproduction in five steps:

\begin{enumerate}
  \item \textbf{CPU audit:} Run \texttt{analyze\_cpu\_distribution.py} in
        \texttt{phase1\_homogeneity/} to reproduce the homogeneity summary.
  \item \textbf{Phase~2 statistics:} Run \texttt{all\_instances\_phase2\_extended.py}
        and compare output against \texttt{experimental\_data/aggregated\_*.csv}.
  \item \textbf{Phase~3 x86:} Run \texttt{phase3\_unified\_opt.py} for one
        configuration and compare against \texttt{averages\_n30.csv}.
  \item \textbf{Phase~3 ARM:} Run \texttt{rerun\_cpu\_verified.py} for one ARM
        configuration and compare against
        \texttt{m3\_graviton\_results\_consolidated.csv}.
  \item \textbf{Figures:} Run \texttt{regenerate\_figures.py} to reproduce
        both paper figures from the embedded LaTeX table data.
\end{enumerate}

\subsection{Dataset Summary}

\begin{table}[!htbp]
\centering
\caption{Benchmark execution counts by experimental phase.}
\label{tab:app_dataset}
\small
\begin{tabular}{lr}
\toprule
\textbf{Phase} & \textbf{Runs} \\
\midrule
Phase~1: CPU verification records & $\approx$4,100 \\
Phase~2: x86 single-instance benchmarks & 5,400 \\
Phase~3 x86: scaling strategy experiments & 540 \\
Phase~3 ARM: Graviton serial and horizontal & 234$^{\dagger}$ \\
Phase~3 ARM: independent replication (CPU-enforced) & 240 \\
Phase~4: chunk size sensitivity & 270 \\
\texttt{sa-east-1} regional replication & 90 \\
\midrule
\textbf{Total timed benchmark executions} & $\approx$\textbf{6,774} \\
\bottomrule
\multicolumn{2}{l}{\footnotesize $^{\dagger}$Includes corrected run counts per Table~III footnotes.} \\
\end{tabular}
\end{table}

\end{document}
